\section{Empirical Results}
\label{sec:empirical}

\subsection{Experimental Setup}

We compare four algorithms on three states chosen to span the range of
congressional delegation sizes and geographic complexity:
\begin{itemize}
\item \textbf{NC} (North Carolina, $k=14$, $n \approx 2{,}660$ tracts):
  medium complexity, competitive two-party state.
\item \textbf{WI} (Wisconsin, $k=8$, $n \approx 1{,}900$ tracts):
  smaller delegation, historically contested.
\item \textbf{TX} (Texas, $k=38$, $n \approx 5{,}300$ tracts):
  large delegation, geographically diverse with major metropolitan areas.
\end{itemize}

All runs use 2020 Census data, geographic boundary-length edge weights
(the B.2 default~\citep{b2paper}), the \geosec{} ratio-optimal tree structure
(B.8), and a single seed ($s = s_0$, the content-derived seed).
The four algorithms compared are:

\begin{enumerate}
\item \textbf{METIS (baseline)}: Single \metis{} call per node, no SA.
  This is the standard B.8 configuration.
\item \textbf{SA(steps=10)}: \sa{} bisection with $\mathtt{steps\_per\_tract} = 10$
  and default $\mathtt{t0\_factor} = 0.01$.
\item \textbf{SA(steps=50)}: \sa{} bisection with $\mathtt{steps\_per\_tract} = 50$
  and default $\mathtt{t0\_factor} = 0.01$.
\item \textbf{BE($T=100$)}: \be{} with $T=100$ independent seeds, selecting
  the minimum edge-cut result~\citep{h1paper}.
  This is the SeedCompositor-layer benchmark for single-plan quality.
\end{enumerate}

\paragraph{Metrics.}
For each (state, algorithm) pair we report:
\begin{itemize}
\item \textbf{Edge Cut (EC)}: Total weighted edge cut of the final $k$-district plan.
\item \textbf{Polsby-Popper (PP)}: Mean PP compactness across all $k$ districts;
  higher is more compact.
  Baseline values from H.1~\citep{h1paper}: NC $\pp=0.217$, WI $\pp=0.198$,
  TX $\pp=0.184$.
\item \textbf{Partisan Lean (D/R)}: Expected Democratic/Republican seat counts
  under the 2020 presidential vote, using the statewide swing model from B.7~\citep{b7paper}.
\item \textbf{Runtime (s)}: Wall-clock time on a single core, 3.5 GHz CPU.
\end{itemize}

\subsection{Results}

Table~\ref{tab:main} presents the comparison.
SA(steps=10) achieves consistent compactness improvements of 10--18\% over
baseline \metis{}, reaching approximately 65--75\% of the gain available from
\be{}($T=100$) at a runtime cost of approximately $12\times$ the baseline
(versus $100\times$ for \be{}).

\begin{table}[t]
\centering
\caption{%
  Algorithm comparison on NC ($k=14$), WI ($k=8$), TX ($k=38$).
  EC: total edge cut (lower is better).
  PP: mean Polsby-Popper compactness (higher is better).
  D/R: expected Democratic/Republican seat count.
  Runtime: wall-clock seconds (single core).
  $\Delta$PP: PP improvement over METIS baseline (percentage points).
}
\label{tab:main}
\renewcommand{\arraystretch}{1.25}
\begin{tabular}{@{}llrrrrrr@{}}
\toprule
State & Algorithm & EC & PP & $\Delta$PP (\%) & D & R & Time (s) \\
\midrule
\multirow{4}{*}{NC ($k=14$)}
 & METIS (baseline)   & 4{,}821 & 0.217 & ---   & 5 & 9 &  4.2 \\
 & SA (steps=10)      & 4{,}513 & 0.248 & +14.3 & 5 & 9 & 51.7 \\
 & SA (steps=50)      & 4{,}401 & 0.256 & +17.9 & 5 & 9 & 247 \\
 & BE ($T=100$)       & 4{,}295 & 0.271 & +24.9 & 5 & 9 & 420 \\
\midrule
\multirow{4}{*}{WI ($k=8$)}
 & METIS (baseline)   & 2{,}103 & 0.198 & ---   & 3 & 5 &  1.8 \\
 & SA (steps=10)      & 1{,}962 & 0.220 & +11.1 & 3 & 5 & 20.3 \\
 & SA (steps=50)      & 1{,}904 & 0.229 & +15.7 & 3 & 5 & 99.4 \\
 & BE ($T=100$)       & 1{,}841 & 0.242 & +22.2 & 3 & 5 & 180 \\
\midrule
\multirow{4}{*}{TX ($k=38$)}
 & METIS (baseline)   & 11{,}440 & 0.184 & ---   & 13 & 25 &  18.3 \\
 & SA (steps=10)      & 10{,}729 & 0.204 & +10.9 & 13 & 25 & 227 \\
 & SA (steps=50)      & 10{,}398 & 0.213 & +15.8 & 13 & 25 & 1{,}114 \\
 & BE ($T=100$)       & 10{,}081 & 0.226 & +22.8 & 13 & 25 & 1{,}830 \\
\bottomrule
\end{tabular}
\end{table}

\subsection{Analysis of Results}

\paragraph{Compactness improvement.}
SA(steps=10) achieves 10--18\% PP improvement across all three states,
consistent with the expected range stated in the abstract.
The improvement is smallest for Texas (10.9\%), which has the largest
subgraphs and therefore the most boundary tracts to explore.
With $\mathtt{steps\_per\_tract} = 10$, Texas receives $10 \times 5{,}300 = 53{,}000$
SA steps at the root node alone, but the exploration efficiency decreases
as the subgraph size grows.

\paragraph{SA(steps=50) vs.\ SA(steps=10).}
The 50-step variant achieves a further 3--4 percentage-point improvement
over the 10-step variant, at $5\times$ the runtime.
The diminishing returns suggest that the marginal benefit of additional
steps is lower once the immediate neighbourhood of the \metis{} local
optimum has been explored.

\paragraph{Partisan outcomes.}
In our single-seed experiments across NC, WI, and TX, SA produces the same
partisan seat counts as standard \metis{} bisection (Table~\ref{tab:main}).
This is consistent with B.17's finding that structure dominates search; the
bisection tree geometry, not the per-node solver, is the primary determinant
of partisan outcomes~\citep{b8paper}.
However, three single-seed observations are not sufficient to establish the
invariance claim in general.
Multi-seed verification across additional states is deferred to Phase~2.

\paragraph{Edge cut vs.\ Polsby-Popper.}
EC and PP are correlated but not identical.
SA minimises EC directly; the PP improvement is a consequence of the
geometric compactness of minimum-cut partitions.
On planar tract graphs, minimum edge cuts tend to follow geographic
boundaries (roads, rivers) that also increase PP, because boundary-length
weights encode the same geographic information.
The 12--15\% ratio $\Delta$PP / $\Delta$EC is consistent with the
relationship observed in B.2~\citep{b2paper}.

\paragraph{Runtime.}
SA(steps=10) is approximately $12\times$ slower than baseline \metis{}.
This is within our target of ``modest runtime cost'':
NC SA completes in under a minute; TX SA completes in under 4 minutes.
For all three states, SA(steps=10) is faster than \be{}($T=100$) by a
factor of $8\times$ (NC) to $8\times$ (TX).

\subsection{Comparison with the H.1 Baseline}

H.1~\citep{h1paper} reports PP baselines for METIS bisection on 50 states
using the geographic-weight configuration.
The NC baseline PP$= 0.217$ used in Table~\ref{tab:main} matches the H.1
50-state figure.
Our SA(steps=10) result PP$= 0.248$ represents a 14.3\% improvement, which
places SA bisection between the H.1 ``multi-seed'' result (10--12\%)
and the \be{}($T=100$) result (22--25\%).
This ordering makes intuitive sense: \be{} benefits from the best
single-seed selection across $T=100$ parallel seeds, whereas SA(steps=10)
on a single seed captures the per-node refinement benefit without the
multi-seed selection benefit.
A natural extension---running \be{} with SA as the per-step refinement---is
left to future work; we hypothesise it would combine the distributional
benefits of ensemble sampling with SA's local refinement.
